problem:  
Let \( (M^n,g) \) be a complete, noncompact Riemannian manifold such that  

1. The sectional curvature satisfies \( \mathrm{sec}_M \ge 0 \) outside a compact subset \( K \subset M \);  
2. There exists a smooth proper function \( f:M \to \mathbb{R} \) with **asymptotically soliton-like behavior**, namely  
   \[
   \big|\mathrm{Ric}_g + \nabla^2 f - \tfrac{1}{2}g\big| = o(r^{-2}),
   \]
   where \(r(x) = d(p,x)\) for some basepoint \(p\);  
3. Each pointed rescaling \( (M, r_i^{-2} g, p) \) with \( r_i\to\infty \) converges in the pointed Gromov–Hausdorff sense to a metric cone \( C(\Sigma) \), with \(\Sigma\) a smooth manifold of dimension \(n-1\).  

**Question:**  
Define \( S := \mathrm{argmin}(f)\), the compact submanifold where the potential \(f\) attains its minimum.  
Does \(S\) act as a **generalized asymptotic soul** in the sense that:  

- \( M \) is diffeomorphic to the total space of the normal bundle \(\nu(S)\); and  
- the diffeomorphism type of \( S \) depends only on the link \(\Sigma\) of the tangent cone at infinity, i.e. the geometry of \( (\Sigma,g_\Sigma) \) uniquely determines the topology of \(S\)?  

In other words, when a manifold is asymptotically Ricci soliton-like and asymptotically conical with curvature nonnegative outside a compact set, does the interaction between \(f\) and the cone at infinity encode a “soul” that controls the global topology of \(M\)?  

Why is it a "good" problem:  
This question proposes a precise, yet uncharted, **analytic analogue of the Cheeger–Gromoll Soul Theorem** in the context where curvature is nonnegative only outside compact sets but is governed asymptotically by a Ricci soliton potential.  Here the potential \(f\) replaces the Busemann function of the classical theory; its level sets provide a canonical exhaustion, and its critical set offers a natural topological candidate for the manifold’s core.  

If an analogue of the Soul Theorem held in this setting, it would unveil a deep connection among three major frameworks:  

- **Topological convexity** from nonnegative curvature (as in the classical Soul Theorem),  
- **Geometric analysis and Ricci flow**, through the asymptotic soliton condition, and  
- **Asymptotic metric geometry**, through tangent cones at infinity.  

Moreover, the question is simple to state in geometric terms but potentially powerful in consequence—it asks whether asymptotic analytic structure rigidifies global topology. A positive resolution would unify analytic and topological descriptions of noncompact manifolds under curvature bounds, extending the reach of both the Soul Theorem and Ricci soliton theory into new terrain.